% Same unit vectors, angular frequency and position pairs as the preceding prose.
\par\noindent\begin{minipage}{\linewidth}
\tagstructbegin{tag=Figure,alttext={沿用正文的二维 RoPE 算例，内容向量均为一、零，角频率为三分之派。左图 Query 在位置零，旋转后为一、零；Key 在位置一，旋转后为二分之一、二分之根号三。右图两位置同时后移两位，Query 在位置二，旋转后为负二分之一、二分之根号三；Key 在位置三，旋转后为负一、零。两图的相对夹角均为六十度，相对位置差均为一，点积均为二分之一。}}
\begin{center}
\begin{tikzpicture}[>=Latex,x=14mm,y=14mm]
\foreach \shift/\title in {0/{位置 $m=0,\ n=1$},8.5/{位置 $m=2,\ n=3$}}{
 \begin{scope}[xshift=\shift cm]
 \draw[->,linegray] (-1.5,0)--(1.65,0);
 \draw[->,linegray] (0,-1.2)--(0,1.7);
 \draw[linegray,dashed] (0,0) circle (1);
 \node[font=\small\sffamily\bfseries] at (0,2.05) {\title};
 \node[font=\small] at (0,-1.6) {$n-m=1$，点积 $1/2$};
 \end{scope}
}
\begin{scope}
 \draw[->,very thick,blue] (0,0)--(0:1) node[below right,font=\small] {Q $(1,0)$};
 \draw[->,very thick,teal] (0,0)--(60:1) node[above,font=\small] {K $(\tfrac12,\tfrac{\sqrt3}2)$};
 \draw[->,flowcolor,thick] (.4,0) arc (0:60:.4);
 \node[font=\small] at (.65,.3) {$60^\circ$};
\end{scope}
\begin{scope}[xshift=8.5cm]
 \draw[->,very thick,blue] (0,0)--(120:1) node[above,font=\small] {Q $(-\tfrac12,\tfrac{\sqrt3}2)$};
 \draw[->,very thick,teal] (0,0)--(180:1) node[below left,font=\small] {K $(-1,0)$};
 \draw[->,flowcolor,thick] (120:.4) arc (120:180:.4);
 \node[font=\small] at (-.65,.3) {$60^\circ$};
\end{scope}
\node[font=\small,text=muted] at (3.0357,-2.35)
 {共同后移两位：两向量都多转 $120^\circ$，相对夹角仍为 $60^\circ$};
\end{tikzpicture}
\captionof{figure}{沿用正文的 $\theta=\pi/3$：$(m,n)$ 从 $(0,1)$ 平移到 $(2,3)$，点积仍为 $1/2$}
\label{fig:rope-relative-rotation}
\end{center}
\tagstructend
\end{minipage}\par
